
Multimodal subsurface velocity inversion combines complementary observations: post-stack migrated images and horizons reveal structural geometry, whereas RMS velocity and well logs calibrate smooth background velocity and absolute scale. Existing methods fuse these cues with one full-field predictor, treating predictable background velocity and structural variation as one task. We identify this mismatch as an asymmetry in conditional complexity and propose Physics-Decoupled Background-Guided Residual Flow Matching (PD-BG-RFM). Our method uses physics-decoupled contrastive learning to construct separate conditions for background-velocity prediction and structural guidance. The background branch deterministically predicts background velocity and reuses it throughout the pipeline: this prediction defines the residual target during training, conditions the residual transport vector field jointly with the structural condition, and composes the final velocity estimate at inference. We derive an exact composition-consistency identity and prove a necessary-and-sufficient condition for residual transport to have lower straight-path target energy than full-field transport. On the globally held-out OpenFWI test set, PD-BG-RFM achieves the best RMSE, SSIM, and high-frequency MAE among matched-input baselines using only 10.6M parameters. Controlled ablations and branch-specific diagnostics support a broader principle for multimodal inverse problems: predict well-constrained components directly and reserve generative transport for structure-bearing variation unexplained by the predicted background.
